\SetPractica{Propiedades y reconocimiento de alcoholes}{Propiedades y reconocimiento de alcoholes}

\section{Introducción teórica}

\section{Objetivos}

\begin{itemize}
    \item Observar algunas propiedades físicas de los alcoholes, así como también realizar reacciones de reconocimiento de alcoholes.
\end{itemize}

\section{Materiales, equipos y reactivos}

\begin{table}[H]
\centering{\small\setstretch{1.4}
\begin{tabular}{|m{0.325\linewidth}|m{0.325\linewidth}|m{0.325\linewidth}|} \hline
    
    \thead{Equipos} & \thead{Materiales} & \thead{Reactivos} \\ \hline
    
    {Balanza} 
    
    &
    
    {Tubos de ensayo medianos \par}
    {Tubos de ensayo grandes \par}
    {Pipetas Pasteur \par}
    {Cápsulas de porcelana \par}
    {Mechero Bunsen \par}
    {Soporte para mechero \par}
    {Malla de asbesto \par}
    {Micropipeta 1mL \par}
    {Pipeta graduada \par}
    {Varilla de vidrio \par}
    {Vaso de precipitado}
    
    & 
    
    {Ácido sulfúrico concentrado \par}
    {Agua destilada \par}
    {Alcohol amílico \par}
    {Alcohol isoamílico \par}
    {Alcohol isopropílico \par}
    {Alcohol terbutílico \par}
    {Benceno \par}
    {Bisulfuro de carbono \par}
    {Etanol absoluto \par}
    {Eter etílico \par}
    {Hidróxido de potasio (\ce{KOH}) \par}
    {Metanol anhídro \par}
    {Reactivo de Lucas \par}
    {Dicromato de potasio (\ce{X}) al 5\% \par}
    {Tetracloruro de carbono}
    
    \\ \hline
\end{tabular}}\end{table}

\section{Procedimientos}

\subsection{Solubilidad de los alcoholes}

\begin{enumerate}
    \item Marque 16 tubos de ensayo medianos con las letras: A1, A2, A3, A4, B1, B2, B3, B4, C1, C2, C3, C4, D1, D2, D3 y D4.
    \begin{itemize}
        \item En los tubos A vierta 0.5 ml de metanol
        \item En los tubos B coloque 0.5 ml de etanol
        \item En los tubos C introduzca 0.5 ml de alcohol isoamílico
        \item En los tubos D vierta 0.5 ml de alcohol amílico
    \end{itemize}
    
    \item Perciba el olor de cada uno de ellos, abanicando con la mano los vapores hacia su nariz
    \begin{itemize}
        \item En los tubos 1 agregue 1 ml de agua destilada
        \item En los tubos 2 añada 1 ml de benceno
        \item En los tubos 3 vierta 1 ml de acetona
        \item En los tubos 4 agregue 1 ml de tetracloruro de carbono.
    \end{itemize}
    
    \item Tape los tubos, agite y déjelos reposar.
\end{enumerate}